\section{Tangent coefficient divisors}
\label{sec:tangent-v171}
We now apply the conductor identity to a coefficient change which is
not diagonal in the given parameters. Throughout this section $k$ is
algebraically closed of characteristic zero. Let $a\geq2$, and let
$T_a\to\Spec k[b,c]$ be the Hilbert graph of
\[
 [s_0:t_0]\longmapsto[t_0^a:b s_0^a:cs_0t_0^{a-1}].
\]
The construction and proof of the chain model are given in
Section~\ref{sec:chain-v171}; in particular, its exceptional divisors
$E_i$ have valuations $v_{E_i}(b)=i$, $v_{E_i}(c)=1$, for $1\leq i\leq a$.
For integers $r,s\geq2$ put
\begin{equation}\label{eq:tangent-cover-v171}
 b=u^r,\qquad c=u+v^s.
\end{equation}
Set
\[
 Y_a=T_a\times_{\Spec k[b,c]}\Spec k[u,v],\qquad Z_a=Y_a^\nu.
\]
The two reduced divisors $u=0$ and $u+v^s=0$ have intersection
multiplicity $s$ at the origin. Thus this is not a change by units of
two transverse boundary coordinates. We do not replace $v$ by an
$s$th root of $c-u$ on the original coefficient surface.

\begin{theorem}[The conductor for tangent covers]
\label{thm:tangent-conductor-v171}
The morphism \eqref{eq:tangent-cover-v171} is finite flat of degree
$rs$, with Jacobian ideal $(u^{r-1}v^{s-1})$. The graph $Y_a$ is
integral and a local complete intersection. For every $i$ there is
exactly one prime divisor $D_i\subset Z_a$ over $E_i$. Its ramification
index $e_i$, residue degree $f_i$, and valuations are as follows:
\begin{align}
 i\leq r:\quad g_i&=\gcd(rs,i),&
 (v_{D_i}(u),v_{D_i}(v),e_i,f_i)
   &=\left(\frac{si}{g_i},\frac{i}{g_i},\frac{rs}{g_i},g_i\right),
 \label{eq:tangent-low-v171}\\
 i\geq r:\quad g_i&=\gcd(r,is),&
 (v_{D_i}(u),v_{D_i}(v),e_i,f_i)
   &=\left(\frac{is}{g_i},\frac{r}{g_i},\frac{rs}{g_i},g_i\right).
 \label{eq:tangent-high-v171}
\end{align}
The two formulas agree at $i=r$. The full conductor on $Z_a$ is
\begin{equation}\label{eq:tangent-conductor-v171}
 \mathfrak c_{Z_a/Y_a}=\OO_{Z_a}\left(-\sum_{i=1}^a k_iD_i\right),
 \qquad
 k_i=\begin{cases}
 \big((rs-1)i-rs\big)/\gcd(rs,i)+1,&i\leq r,\\
 \big((r-1)is-r\big)/\gcd(r,is)+1,&i\geq r.
 \end{cases}
\end{equation}
No other conductor component or additional condition in codimension
two occurs. The generic normalization defect over $E_i$ is $f_i k_i/2$.
These assertions hold for all $a,r,s$ in the stated range, not only for
monomial arcs in the coefficient surface.
\end{theorem}
\begin{proof}
The intermediate extension defined by $u^r=b$, followed by
$v^s=c-u$, is free with basis $u^jv^l$, $0\leq j<r$, $0\leq l<s$.
The total algebra is $k[u,v]$, and its fraction field has degree $rs$
over $k(b,c)$. Differentiating gives the determinant
$rsu^{r-1}v^{s-1}$. The graph and local complete-intersection claims
follow from Theorem~\ref{thm:duality-conductor-v171} and the smoothness
of $T_a$.

Fix $i$, write $w=b/c^i$, and localize at the generic point of $E_i$.
The residue field is $k(w)$, with $w$ transcendental. At a prime of the
normalization above this DVR let $\alpha=v(u)$, $\beta=v(v)$, and
$e=v(c)$. Then
\begin{equation}\label{eq:tangent-valuation-relations-v171}
 r\alpha=ie,\qquad v(u+v^s)=e.
\end{equation}
The fundamental equality for the finite normalization of a DVR is
$\sum ef=rs$. We shall obtain lower bounds on a single $e$ and $f$
whose product is already $rs$; this also proves uniqueness of the
prime, rather than assuming it from a Newton polygon.

Suppose first that $i<r$. Then $\alpha<e$, so cancellation in $u+v^s$
forces $s\beta=\alpha$ and $\overline{u/v^s}=-1$. If $g=\gcd(rs,i)$,
the valuation equalities force $e$ to be a positive multiple of $rs/g$.
The unit character $v^{rs/g}/c^{i/g}$ has residue $\eta$ with
$\eta^g=(-1)^r w$. The polynomial $T^g-(-1)^r w$ is irreducible over
$k(w)$, for example by the $w$-adic Eisenstein criterion. Hence $f\geq g$.
The fundamental equality now forces $e=rs/g$, $f=g$, and a single
prime above $E_i$. This proves \eqref{eq:tangent-low-v171} for $i<r$.

If $i=r$, the residue $\zeta=\overline{u/c}$ satisfies $\zeta^r=w$.
It has degree $r$ and is not $1$. Since
$v^s=c(1-u/c)$, one has $s\beta=e$. Consequently $e\geq s$ and
$f\geq r$, and again equality and uniqueness follow. Both displayed
formulas give $(\alpha,\beta,e,f)=(s,1,s,r)$.

If $i>r$, then $\alpha>e$, so $s\beta=e$ and
$\overline{c/v^s}=1$. Set $g=\gcd(r,is)$. The relations force
$e\geq rs/g$, and the residue of $u^{r/g}/v^{is/g}$ has $g$th power
$w$. The same irreducibility and degree argument gives
$e=rs/g$, $f=g$ and \eqref{eq:tangent-high-v171}.

The conductor coefficient is now
$(r-1)\alpha+(s-1)\beta-e+1$ by
Theorem~\ref{thm:duality-conductor-v171}. Substitution yields
\eqref{eq:tangent-conductor-v171}. The theorem is an equality of ideals,
so no extra intersection condition is being inferred merely from
these orders. Away from the exceptional locus the pullback is the
normal plane $\Spec k[u,v]$. The defect formula follows from
Corollary~\ref{cor:defect-cycle-v171}.
\end{proof}

The residue degrees in this theorem describe extensions of function
fields of divisors. They do not count prime divisors. Over a general
nonzero closed residue, algebraic closedness and characteristic zero
turn the indicated power maps into $g_i$ distinct closed lifts. Special
points require the local models, which we compute next in an all-order
nonmonomial case.

\subsection{The complete square-tangency model}
For the rest of this section take $r=s=2$. We describe the full
normal surface and its scheme fibre, not just its divisorial valuations.

\begin{theorem}[Square tangency in every contact order]
\label{thm:square-tangency-v171}
For $b=u^2$, $c=u+v^2$, the surface $Z_2$ is covered by
\begin{align}
 B_0&=k[u,v,k_0]/(u+v^2-u^2k_0),\label{eq:square-chart-zero-v171}\\
 B_1&=k[e,z,v]/(v^2-ez(z-1)),\label{eq:square-chart-one-v171}\\
 B_2&=k[c,q,v]/(v^2-c(1-q)).\label{eq:square-chart-two-v171}
\end{align}
The maps to $\Spec k[u,v]$ are the evident one on $B_0$,
$u=ez$ on $B_1$, and $u=cq$ on $B_2$. The overlaps are
\[
 k_0=e^{-1},\quad z=uk_0;
 \qquad c=ez^2,\quad q=z^{-1}.
\]
There are precisely two singular points, both of type $A_1$; in the
middle chart they are $(e,z,v)=(0,0,0)$ and $(0,1,0)$.
The conductors on these charts are respectively
\begin{equation}\label{eq:square-chart-conductors-v171}
 B_0,\qquad eB_1,\qquad cB_2.
\end{equation}
For every $a>2$, the normal surface $Z_a$ is obtained from $Z_2$ by
$a-2$ successive ordinary blow-ups at the specified smooth endpoint
of the exceptional chain. Its fibre over $(u,v)=(0,0)$ is a reduced
nodal chain of $a$ projective lines, with no embedded points. The two
$A_1$ points persist and no new singularities occur. For every $a\geq2$,
\begin{equation}\label{eq:square-global-conductor-v171}
 \mathfrak c_{Z_a/Y_a}
 =\OO_{Z_a}(-2D_2-3D_3-\cdots-aD_a).
\end{equation}
In particular $D_1$ and $D_2$ have the same values
$(v(u),v(v))=(2,1)$ but different values $4$ and $2$ on $u+v^2$.
A fan in the two original monomial weights cannot distinguish them.
\end{theorem}
\begin{proof}
The three charts of $T_2$ in Section~\ref{sec:chain-v171} give raw
pullback rings
\begin{align*}
 A_0&=k[u,v,k_0]/(u+v^2-u^2k_0),\\
 A_1&=k[e,h,u,v]/(u^2-e^2h,\ u+v^2-eh),\\
 A_2&=k[u,v,w]/(u^2-(u+v^2)^2w).
\end{align*}
They cover the entire pullback because they are base changes of an
open cover of the blow-up. Their integrality has already been proved
by finite flatness, so all fractional expressions below lie in their
single fraction field.

The ring $A_0$ is smooth: where $u\ne0$ the derivative with respect to
$k_0$ is nonzero, and where $u=0$ the derivative with respect to $u$
is $1$. For $A_1$ adjoin $z=u/e$, integral because $z^2=h$. Eliminating
$u,h$ gives $B_1$. For $A_2$ put $c=u+v^2$ and adjoin $q=u/c$,
integral because $q^2=w$. This gives $B_2$. These extensions are
finite and birational. The two displayed hypersurfaces are smooth
except at the points specified in the theorem; their isolated
singularities and the hypersurface $S_2$ property prove normality.
At $z=0$ and $z=1$ the factors $z-1$ and $z$, respectively, are units,
so completion and a unit change of $e$ give $v^2=ez$ or $v^2=e(z-1)$.
These are $A_1$ singularities. The point $c=v=0,q=1$ of $B_2$ is the
same second point. The fractional identities on overlaps prove the
stated gluing of the finite normalizations.

We verify the conductors at the intersections directly. Over
$D=k[e,h]$, the normal ring $B_1$ is free with basis
$1,z,v,zv$, using $z^2=h$ and $v^2=e(h-z)$. The subring $A_1$ is
exactly the free sublattice with basis $1,ez,v,ezv$. Thus an element
$A+Bz+Cv+Dzv$ belongs to $A_1$ if and only if $B,D$ are divisible by
$e$. Requiring its product with $z$ also to belong to $A_1$ forces
$A,C$ to be divisible by $e$. These necessary conditions say that
the element is in $eB_1$, and they are sufficient because
$eB_1\subset A_1$ and $eB_1$ is a $B_1$-ideal. This proves the
middle equality in \eqref{eq:square-chart-conductors-v171}.
Over $k[c,w]$, use the basis $1,q,v,qv$ with $q^2=w$ and
$v^2=c(1-q)$. The sublattice is $1,cq,v,cqv$; the identical argument
gives $cB_2$. On the overlap $c/e=z^2$ is a unit, so these are the
same ideal. On the first overlap $e$ is a unit, giving $B_0$.
This is an explicit intersection calculation, independent of any
extension of a generic divisor formula.

The complete fibre is equally explicit:
\begin{equation}\label{eq:square-full-fibre-v171}
 B_0/(u,v)=k[k_0],\qquad
 B_1/(u,v)=k[e,z]/(ez),\qquad
 B_2/(u,v)=k[q].
\end{equation}
For the last equality, $cq=0$ and $c(1-q)=0$ imply $c=0$.
The overlaps compactify the two affine axes of the middle node into
two projective lines. These rings are reduced, pure one-dimensional,
and have no embedded points. Notice that the second singularity of
the ambient surface lies at a smooth point of the fibre.

It remains to prove the assertion for all $a$, rather than infer it
from low-order pictures. The morphism $T_{j+1}\to T_j$ blows up the
last point with coordinates $(c,w_j)=(0,0)$, where $w_j=b/c^j$.
Its inverse image on $Z_j$ is a unique smooth endpoint with local
parameters $v,q_j$. In a neighbourhood of this point put
\begin{equation}\label{eq:tangent-induction-v171}
 H_j=1-v^{j-2}q_j,\qquad
 c=\frac{v^2}{H_j},\quad
 u=c v^{j-2}q_j,\quad
 w_j=q_j^2H_j^{j-2}.
\end{equation}
For $j=2$ these are the last chart above, localized at $q_2=0$;
$H_2=1-q_2$ is a unit. The pullback of $(c,w_j)$ is
$(v^2,q_j^2)$. Its normalized blow-up is the ordinary blow-up of
$(v,q_j)$: the integral closure of $(v^2,q_j^2)^n$ is
$(v,q_j)^{2n}$ for every $n$, by its two-dimensional Newton polygon,
and Proj of this Veronese algebra is the ordinary blow-up.

For completeness this computes the required normalized base change,
not a substitute birational model. The finite morphism $Z_j\to T_j$
is flat. Indeed $Z_j$ is a normal surface and hence Cohen--Macaulay,
whereas $T_j$ is smooth; a system of parameters of the latter is a
regular sequence on the former, giving local freeness over the
regular local ring. Flat base change of the blow-up is therefore the
blow-up of the pulled-back ideal. The finite birational map from this
pullback to $Y_{j+1}$ induces the same function field, so their
normalizations agree by transitivity of integral dependence.

On the last chart of the ordinary blow-up set $q_j=vq_{j+1}$.
Formula \eqref{eq:tangent-induction-v171} becomes the same formula
with $j+1$, proving the induction and uniqueness of the next centre.
The old fibre ideal at that centre is $(u,v)=(v)$, so its pullback is
the reduced total transform of one smooth component. The blow-up adds
one projective line and one node, not an embedded point or a multiple
component. The centre is disjoint from the two $A_1$ points; thus
no singularities are added. Finally Theorem~\ref{thm:tangent-conductor-v171}
gives $k_1=0$ and $k_i=i$ for $i\geq2$, as well as the valuation
comparison for $D_1,D_2$. This proves every assertion.
\end{proof}

\begin{corollary}[One Hilbert limit, different normalization germs]
\label{cor:tangent-lost-labels-v171}
For every $a\geq2$, two coefficient arcs of the square-tangency family
can have the same closed coefficients and the same embedded Hilbert
limit, while their normalization lifts have nonisomorphic completed
local rings: one is regular and the other has an $A_1$ singularity.
Both arcs are generically base-point-free as projective maps.
\end{corollary}
\begin{proof}
On $D_2$ in the middle chart the old Hilbert coordinate is $h=z^2$.
The points $z=1$ and $z=-1$ therefore map to the same point $e=0,h=1$
of $T_2$, and hence give the same embedded curve. The former is an
$A_1$ point of $Z_2$, whereas the latter is smooth. They are untouched
by the later endpoint blow-ups.

For the point $z=-1$ take $v=\tau$, $e=\tau^2/2$, $u=-\tau^2/2$.
For the point $z=1$ take
\[
 z=1+\tau,\qquad v=\tau,\qquad
 e=\frac{\tau}{1+\tau},\qquad u=\tau.
\]
Both satisfy $v^2=ez(z-1)$ over $k[[\tau]]$ and have
$b=u^2$, $c=u+v^2$. In both cases $u$ is nonzero over $k((\tau))$.
If the homogeneous source coordinate $s_0$ is nonzero, the second
coordinate $u^2s_0^a$ is nonzero; if $s_0=0$, the first coordinate
$t_0^a$ is nonzero. Thus the entire projective map, including source
infinity, is base-point-free on the generic fibre. The regular and
$A_1$ completed local rings have different embedding dimensions, so
they cannot be isomorphic.
\end{proof}
